Binary rewriting consists in altering an already compiled executable without having access to the source code,
        while ensuring that the modified binary remains executable.
        The process of binary rewriting can be classified into Static Binary Rewriting, if the binary is stored in
        persistent memory while being rewritten, and Dynamic Binary Rewriting, if the binary is being executed.
        The process of binary rewriting can be divided into four main steps \cite{binaryrewriting}:
        \begin{enumerate}
            \item \textit{Parsing}: In this phase the main goal is to retrieve information about the binary.
                In particular, the output of this phase is a raw stream of instructions (in machine code) and data.
                This information is then used in the Analysis step.

            \item \textit{Analysis}: The information gained from the Parsing phase is analyzed to reconstruct the program structure.
            For example, the raw stream of machine code is passed through a disassembler that reconstructs the
            assembly instructions. A structural recovery is performed, which includes function detection and \gls{CFG} reconstruction,
            and labels and symbols are extracted from the raw stream of data.

            \item \textit{Transformation}: Once the \gls{CFG} has been reconstructed and function and data boundaries have been recovered,
                the binary rewriter can modify the binary at the \gls{IP}. These are user-defined locations in the binary where
                the rewriter must add, modify or delete an instruction. The \gls{IP} can be coarse- or fine-grained.
                In the first case, for example, the instrumentation could be inserted at branch level (at the beginning of each \gls{BB}).
                In the second case, instead, any instruction could potentially be modified.

            \item \textit{Code Generation}: In this last step, the changes made in the previous steps must be integrated into the
                binary while preserving its executability. This can be done in different ways. For example, the changes could be
                inserted into a new section of the executable, redirecting execution to the new code at the \gls{IP}.
                Alternatively, this step could also be performed by reassembling a completely new executable starting from the
                changes made at the \gls{IP}. This second approach is much harder, because it requires preserving all the
                references inside the executable and preserving the executable layout.
        \end{enumerate}
        \begin{figure}[ht]
            \caption{An example of a Static Binary Rewriting Workflow}
            \centering
            \includegraphics[width=200pt]{BinaryRewritingScheme.png}
            \label{fig:static_rewrite}
        \end{figure}
        In \cref{fig:static_rewrite}, a general scheme for static binary rewriting can be seen.
        In this case all the phases occur prior to execution, so the cost in terms of overhead
        is paid only once, before the execution. The main problems are disambiguating between code and data
        (i.e. deciding whether a stream of bytes represents machine code instructions or constants placed inline in the .text section)
        and resolving indirect branches \cite{binaryrewriting,andriesse2016}.
        This problem is generally undecidable and can only be solved using heuristics.
        There are three main families of static transformation:
        \begin{itemize}
            \item \textbf{Direct}: Consists in directly modifying the instruction. Works well with RISC systems with instructions having
                the same length. Generally fails when adding or removing instructions, because this does not preserve the offset into the .text section.
            \item \textbf{Minimal-invasive}: The new code is inserted into a new section of the binary, and at the \gls{IP}
                an unconditional jump to the new section is placed.
                The main limitation of this approach arises when the new section is far from the \gls{IP} and a near jump (necessary to reach the new section)
                uses a different number of bytes than the instruction at the instrumentation point.
                In fact, in x86-64 assembly a short jump uses 2 bytes while a near jump uses up to 5 bytes.
                So, if a short jump is needed at the instrumentation point for alignment reasons,
                but the new section is too far away and a near jump is needed instead, it is necessary to add a \textbf{trampoline}.
                This approach can be seen in \cref{fig:trampoline_instrumentation}.
                Trampoline-based instrumentation is one of the most common approaches to binary rewriting \cite{binaryrewriting}.

            \item \textbf{Full-translation}: The binary is lifted into an \gls{IR} that makes it possible to add an \gls{IP} at
                instruction level. The \gls{IR} is then reassembled into machine code after the changes are made.
        \end{itemize}
        \begin{figure}[ht]
            \caption{An example of Static Instrumentation using the trampoline technique}
            \centering
            \includegraphics[width=400pt]{trampolineInstrumentation.png}
            \label{fig:trampoline_instrumentation}
        \end{figure}
        With regard to Dynamic Binary Rewriting, the Analysis and Transformation phases are executed at runtime for each instruction.
        In this case, none of the problems that arise in Static Binary Rewriting occur, since the binary is being executed,
        so indirect branches can be resolved using the values in the registers, and there is no need to
        disambiguate between constants inside the .text section and instructions.
        The main drawback of this approach is the prohibitive overhead caused by the dynamic analysis and transformation of each instruction,
        together with the fact that only the taken branches can be modified.
        In fact, if a branch is not actually taken during execution, its instructions never reach the rewriter.\\
        \\
        A specific type of binary rewriting is Binary Instrumentation.
        Binary instrumentation uses binary rewriting techniques to add instructions that track binary behaviour,
        log events, or, in the case of this thesis, insert edge coverage tracing instructions.